% Biên dịch được bằng pdfLaTeX (mặc định của Overleaf), LuaLaTeX hoặc XeLaTeX.
% Bản nộp ARR phải là tiếng Anh; bản này là bản nháp nội dung tiếng Việt.
\documentclass[11pt]{article}

% "review" khi nộp (ẩn danh + số dòng); "final" cho camera-ready; "preprint" cho bản không ẩn danh.
\usepackage[review]{acl}

\usepackage{iftex}
\ifPDFTeX
  % pdfLaTeX: tiếng Việt cần bảng mã T5 (gói vntex, có sẵn trong TeX Live/Overleaf)
  \usepackage{times}
  \usepackage[T5]{fontenc}
  \usepackage[utf8]{inputenc}
  \usepackage[vietnamese]{babel}
\else
  % LuaLaTeX / XeLaTeX
  \usepackage[bidi=default]{babel}
  \babelprovide[import,main]{vietnamese}
  % font theo tên file (có trong mọi TeX Live), không phụ thuộc font cài trên máy
  \babelfont{rm}[Extension=.otf, UprightFont=*-regular, BoldFont=*-bold, ItalicFont=*-italic,
    BoldItalicFont=*-bolditalic]{texgyretermes} % tương đương Times, có đủ dấu tiếng Việt
  \babelfont{tt}[Extension=.otf, UprightFont=*-regular, BoldFont=*-bold, ItalicFont=*-italic,
    BoldItalicFont=*-bolditalic]{texgyrecursor}
\fi
\usepackage{microtype}

\usepackage{amsmath}
\usepackage{amssymb}
\usepackage{booktabs}
\usepackage{array}
\newcolumntype{L}[1]{>{\raggedright\arraybackslash}p{#1}}
\usepackage{multirow}
\usepackage{graphicx}
\usepackage{xcolor}
\usepackage{tikz}
\usetikzlibrary{positioning,arrows.meta,fit,backgrounds,calc}
% figure palette: categorical slots 1-3 (blue, orange, aqua), an ordinal blue pair, ink and chrome grays
\definecolor{vzblue}{HTML}{2A78D6}
\definecolor{vzorange}{HTML}{EB6834}
\definecolor{vzaqua}{HTML}{1BAF7A}
\definecolor{vzblueL}{HTML}{86B6EF}
\definecolor{vzblueD}{HTML}{1C5CAB}
\definecolor{inkP}{HTML}{0B0B0B}
\definecolor{inkS}{HTML}{52514E}
\definecolor{inkM}{HTML}{898781}
\definecolor{gridc}{HTML}{E1E0D9}
\definecolor{axisc}{HTML}{C3C2B7}
% data figures are written by scripts/paper_figures.py into <EVAL_DIR>/paper/figures/; copy the PDFs into
% paper/figures/ (next to the fig_*.tex files). Until then each shows a placeholder box.
\newcommand{\datafigure}[2]{\IfFileExists{#1}{\includegraphics[width=#2]{#1}}{%
  \fbox{\parbox[c][2.2cm][c]{0.9#2}{\centering\small\textcolor{red}{[TBD]} \texttt{#1}\\
  (sinh bởi \texttt{scripts/paper\_figures.py} sau lần chạy đầy đủ)}}}}

\newcommand{\arm}[1]{\texttt{#1}}
\newcommand{\tbd}{\textcolor{red}{\textbf{[TBD]}}}
\newcommand{\todo}[1]{\textcolor{red}{[TODO: #1]}}
\newcommand{\R}{\mathbb{R}}

% Working title (EN): Amortized Utility Attribution for Reader-Agnostic Context Compression
\title{Chưng cất mức đóng góp của từng đoạn văn để nén ngữ cảnh\\
không phụ thuộc mô hình đọc cho hỏi đáp tiếng Việt và đa ngôn ngữ}

\author{Tác giả ẩn danh \\
  Đơn vị ẩn danh \\
  \texttt{email@domain}}

\begin{document}
\maketitle

% ============================================================================
\begin{abstract}
Nén ngữ cảnh cho hệ thống sinh tăng cường truy xuất (RAG) thường được huấn luyện
trên các tín hiệu thay thế như độ liên quan truy vấn--đoạn văn, độ bất ngờ của
token, hoặc nhãn ``đoạn chứa đáp án''. Các tín hiệu này không đo trực tiếp điều
ta quan tâm: một đoạn văn giúp mô hình đọc (reader) trả lời đúng đến mức nào. Các
phương pháp truy vết đóng góp của ngữ cảnh (context attribution) như ContextCite đo được đại lượng này bằng cách che
ngẫu nhiên ngữ cảnh và khớp một mô hình thay thế tuyến tính, nhưng mỗi ví dụ
cần hàng chục lần gọi reader. Vì vậy chúng không dùng được làm bộ nén lúc suy
luận. Chúng tôi kiểm định giả thuyết rằng \emph{mức đóng góp của từng đoạn} (utility), tức hệ số hồi quy Ridge
của F1 đầu ra theo mặt nạ đoạn ngẫu nhiên, có thể được \emph{chưng cất}
(amortize: trả chi phí đo một lần lúc huấn luyện thay vì đo lại mỗi lần dùng) vào một bộ tỉa cross-encoder chỉ chạy một lượt, và rằng nhãn tổng
hợp từ nhiều reader cho bộ nén ít phụ thuộc reader hơn. Phương pháp gồm hai giai
đoạn: (A) sinh nhãn mức đóng góp bằng $K\!\approx\!64$ mặt nạ trên mỗi tài liệu với ba
reader, và (B) chưng cất các nhãn này vào một bộ tỉa dựa trên
\texttt{bge-reranker-v2-m3} bằng hàm mất mát ListNet. Chúng tôi đánh giá trên
năm nguồn dữ liệu tiếng Việt và tiếng Anh, gồm cả một câu hỏi và nhiều bước suy luận, ở ngân sách
$1/4$ và $1/8$ token, với bốn reader trong đó có một reader không tham gia sinh
nhãn, và so sánh với bảy bộ nén đã công bố (Provence, XProvence, RECOMP, EXIT, LLMLingua,
LongLLMLingua, LLMLingua-2) ở cùng ngân sách. Các giả thuyết được đăng ký trước dưới dạng họ kiểm định bootstrap cụm
ghép cặp có hiệu chỉnh Holm. \tbd{} Kết quả chính.
\end{abstract}

% ============================================================================
\section{Giới thiệu}
\label{sec:intro}

Các hệ thống RAG và tác tử dựa trên mô hình ngôn ngữ lớn (LLM) thường đưa vào
reader hàng chục đoạn văn, trong khi chỉ một phần nhỏ thực sự cần cho câu trả
lời. Nén ngữ cảnh, tức giữ lại một tập con của ngữ cảnh trong một ngân sách
token, giúp giảm chi phí suy luận và giảm nhiễu \citep{liu2024lost}. Câu hỏi then
chốt là \emph{bộ nén nên được huấn luyện để giữ cái gì}.

Phần lớn các bộ nén hiện nay tối ưu một \emph{tín hiệu thay thế}. Họ LLMLingua
dùng độ bất ngờ (perplexity) hoặc nhãn chưng cất từ LLM ở mức token
\citep{jiang2023llmlingua,jiang2024longllmlingua,pan2024llmlingua2}. Provence và
XProvence dùng độ liên quan của reranker cùng nhãn gán chuỗi
\citep{chirkova2025provence,xprovence2026}. EXIT và RECOMP dùng nhãn câu hỗ trợ
hoặc câu chứa đáp án \citep{hwang2025exit,xu2024recomp}. Các tín hiệu này tương
quan với điều ta muốn nhưng không trùng với nó. Một đoạn ``liên quan'' có thể
thừa, và một đoạn cầu nối (bridge) trong câu hỏi nhiều bước có thể không chứa đáp án
nhưng vẫn bắt buộc phải có (Hình~\ref{fig:example}).

\begin{figure}[t]
\centering
% Figure 2: why answer-span labels and relevance are proxies (illustrative values, not measurements).
% Emphasis form: the utility column is the point (blue), the two proxies are context (gray / marks).
\begin{tikzpicture}[
    font=\scriptsize,
    para/.style={anchor=north west, text width=3.05cm, align=left, inner sep=0pt, text=inkP},
    head/.style={font=\scriptsize\bfseries, text=inkS, align=center, anchor=south},
  ]
  \def\rowh{9.5mm}
  \def\barh{2.4mm}
  \def\barmax{1.05cm}
  % column x positions
  \def\xspan{3.75cm}
  \def\xrel{4.35cm}
  \def\xbeta{5.95cm}

  \node[anchor=south west, text width=7.3cm, align=left, inner sep=0pt, font=\footnotesize, text=inkP]
    at (0, 0.8cm) {\textbf{Câu hỏi:} Đạo diễn của phim \emph{Mùi đu đủ xanh} sinh năm nào?\quad
    \textbf{Đáp án:} 1962};
  \node[head] at (\xspan, 0.05cm) {Chứa\\ đáp án};
  \node[head] at ($(\xrel,0.05cm)+(0.5*\barmax,0)$) {Độ liên quan\\ (reranker)};
  \node[head] at ($(\xbeta,0.05cm)+(0.5*\barmax,0)$) {Mức đóng góp\\ $\beta$ (của chúng tôi)};
  \draw[axisc, line width=0.4pt] (0,0) -- (7.3cm,0);

  % rows: text / contains answer / relevance / beta
  \foreach \i/\txt/\hasans/\rel/\bval in {
      1/{\textbf{P1} (cầu nối). \emph{Mùi đu đủ xanh} (1993) là phim của đạo diễn Trần Anh Hùng.}/0/0.85/0.90,
      2/{\textbf{P2}. Trần Anh Hùng sinh năm 1962, là đạo diễn người Pháp gốc Việt.}/1/0.40/0.95,
      3/{\textbf{P3}. \emph{Mùi đu đủ xanh} đoạt giải Caméra d'Or tại Cannes năm 1993.}/0/0.90/0.05,
      4/{\textbf{P4}. Đà Lạt là thành phố thuộc tỉnh Lâm Đồng.}/0/0.05/0.02} {
    \pgfmathsetlengthmacro{\ytop}{-1.5mm-(\i-1)*\rowh}
    \pgfmathsetlengthmacro{\ymid}{\ytop-3.2mm}
    \node[para] at (0, \ytop) {\txt};
    \ifnum\hasans=1
      \node[text=inkP, font=\small] at (\xspan, \ymid) {\checkmark};
    \else
      \node[text=inkM] at (\xspan, \ymid) {--};
    \fi
    % baselines, then bars (square at the baseline, rounded data end)
    \draw[axisc, line width=0.4pt] (\xrel, \ymid-2mm) -- (\xrel, \ymid+2mm);
    \draw[axisc, line width=0.4pt] (\xbeta, \ymid-2mm) -- (\xbeta, \ymid+2mm);
    \fill[inkM, rounded corners=0.7pt] (\xrel, \ymid-0.5*\barh) rectangle ++(\rel*\barmax, \barh);
    \fill[vzblue, rounded corners=0.7pt] (\xbeta, \ymid-0.5*\barh) rectangle ++(\bval*\barmax, \barh);
  }
\end{tikzpicture}

\caption{Ví dụ nhiều bước (giá trị minh họa, không phải số đo). Nhãn ``đoạn chứa đáp án'' chỉ
đánh dấu P2 và bỏ sót đoạn cầu nối P1; độ liên quan xếp P3 cao dù P3 không giúp trả lời. Mức
đóng góp $\beta$ đo điều reader thực sự cần: cả P1 lẫn P2.}
\label{fig:example}
\end{figure}

Đại lượng ta thực sự muốn là \emph{mức đóng góp} (utility) của từng đoạn: nếu bỏ đoạn $c_i$ thì chất lượng
câu trả lời của reader giảm bao nhiêu. ContextCite \citep{cohenwang2024contextcite}
cho thấy có thể ước lượng đại lượng này bằng cách che ngẫu nhiên ngữ cảnh rồi khớp
một mô hình thay thế tuyến tính thưa. Tuy nhiên, mỗi ví dụ cần 32--256 lần gọi
reader, nên đây là công cụ \emph{giải thích}, không phải bộ nén. AT2
\citep{cohenwang2025at2} chưng cất phép truy vết này vào trọng số các đầu chú ý của
chính reader và nhắm tới log-xác suất của câu trả lời. LooComp dùng bộ tỉa
cross-encoder riêng nhưng học từ nhãn nhị phân leave-one-out \citep{looComp2026}.
Theo hiểu biết của chúng tôi, chưa có công trình nào \emph{chưng cất hệ số
surrogate của F1 đầu ra, thu được từ mặt nạ ngẫu nhiên, vào một bộ tỉa độc lập
với reader}. Chúng tôi gọi quá trình này là \emph{chưng cất mức đóng góp}: chi phí
đo (hàng chục lần gọi reader cho mỗi tài liệu) chỉ phải trả một lần trên dữ liệu
huấn luyện, còn lúc sử dụng bộ tỉa chỉ cần chạy một lượt.

Vấn đề thứ hai là \emph{phụ thuộc reader}. Nhãn mức đóng góp được đo trên một reader
cụ thể. \citet{readerScaling2026} chỉ ra rằng một bộ nén cố định có thể che đi
tới 80\% mức cải thiện khi nâng cấp reader và đảo 31\% thứ hạng mô hình. Nếu bộ
nén học ``khẩu vị'' của một reader, nó có thể lãng phí năng lực của reader mạnh
hơn.

\paragraph{Giả thuyết trung tâm.} Bài báo kiểm định giả thuyết sau (phát biểu
hình thức ở \S\ref{sec:hypotheses}):
\begin{quote}\itshape
Mức đóng góp biên của từng đoạn vào F1 đầu ra, ước lượng bằng surrogate cộng tính
trên mặt nạ ngẫu nhiên, là một tín hiệu giám sát (i) \textbf{học được} từ cặp
(câu hỏi, ngữ cảnh) mà không cần gọi reader lúc suy luận, (ii) \textbf{chứa thông
tin vượt quá} nhãn ``đoạn chứa đáp án'' khi câu hỏi cần nhiều đoạn, và (iii)
\textbf{phần lớn được chia sẻ} giữa các reader, nên gộp nhãn từ nhiều reader cho một
bộ nén bảo toàn tốt hơn lợi ích khi nâng cấp reader.
\end{quote}

\paragraph{Đóng góp.}
\begin{enumerate}
  \item Chúng tôi đề xuất \textbf{chưng cất mức đóng góp của từng đoạn} (amortized utility
  attribution). Phương pháp sinh nhãn mức đóng góp theo đoạn bằng mặt nạ ngẫu nhiên, surrogate Ridge của F1 và chuẩn
  hóa z trong tài liệu, rồi chưng cất các nhãn này vào một bộ tỉa cross-encoder
  một lượt. Chi phí suy luận giảm từ $K{+}1$ lần gọi reader xuống một lượt encoder
  (\S\ref{sec:method}).
  \item Chúng tôi đưa ra \textbf{nhãn tổng hợp nhiều reader} và đánh giá mức bảo
  toàn nâng cấp (upgrade retention) trên một reader \emph{không tham gia} sinh
  nhãn (\S\ref{sec:method-ens}).
  \item Chúng tôi xây dựng một \textbf{giao thức kiểm định đăng ký trước}: ba giả thuyết cho
  ba câu hỏi nghiên cứu, mỗi giả thuyết có quy tắc quyết định và được kiểm định bằng bootstrap cụm
  ghép cặp với hiệu chỉnh Holm; kết quả được báo cáo theo từng nguồn dữ liệu, không gộp (\S\ref{sec:hypotheses}, \S\ref{sec:setup}). Phương pháp được so với \textbf{bảy
  bộ nén đã công bố} ở ba mức chi tiết (đoạn, câu, token), dùng checkpoint gốc của tác giả và
  cùng một ngân sách token (Phụ lục~\ref{app:baselines}).
  \item Chúng tôi đánh giá \textbf{tiếng Việt} ở hai dạng: một bước (haystack xây từ
  UIT-ViQuAD và XQuAD-vi) và nhiều bước (VIMQA). Phương pháp được so với XProvence
  zero-shot và LLMLingua-2 ở cùng ngân sách.
\end{enumerate}

% ============================================================================
\section{Công trình liên quan}
\label{sec:related}

\paragraph{Nén ngữ cảnh.} Nén ở mức token loại bỏ các token có độ bất ngờ thấp
\citep{li2023selectivecontext,jiang2023llmlingua,jiang2024longllmlingua} hoặc học
một bộ phân loại token từ nhãn chưng cất \citep{pan2024llmlingua2}. Nén trích xuất
chọn câu hoặc đoạn: RECOMP \citep{xu2024recomp}, EXIT \citep{hwang2025exit},
Provence \citep{chirkova2025provence} và phiên bản đa ngôn ngữ XProvence
\citep{xprovence2026}, cũng như LooComp \citep{looComp2026}. Nén sinh tóm tắt tài
liệu bằng một mô hình sinh: RECOMP abstractive và CompAct \citep{yoon2024compact}. Một số phương pháp
gần đây tối ưu trực tiếp theo đầu ra. CORE-RAG huấn luyện một bộ nén sinh bằng
học tăng cường với phần thưởng EM \citep{coreRag2026}. ECoRAG dùng tín hiệu
``evidentiality'' \citep{ecorag2025}. PoC dự đoán hiệu năng để chọn tỉ lệ nén
\citep{poc2026}. Nén mềm như gist tokens và ICAE \citep{mu2023gist,ge2024icae}
nằm ngoài phạm vi bài này vì chúng đòi hỏi can thiệp vào reader.

\paragraph{Truy vết đóng góp của ngữ cảnh.} Datamodels \citep{ilyas2022datamodels} và
ContextCite \citep{cohenwang2024contextcite} ước lượng ảnh hưởng của từng thành
phần bằng hồi quy tuyến tính trên các ablation ngẫu nhiên. AT2
\citep{cohenwang2025at2} học tổ hợp các đầu chú ý để tái tạo điểm đóng góp
trong một lượt. Chúng tôi dùng lại cơ chế ablation ngẫu nhiên kết hợp surrogate
tuyến tính, nhưng khác ở ba điểm: (a) đầu ra là \emph{F1 của câu trả lời được
sinh}, không phải log-xác suất; (b) mô hình được chưng cất là một \emph{bộ tỉa riêng}
không truy cập trạng thái của reader; (c) nhãn được gộp từ nhiều reader.

\paragraph{Phụ thuộc reader.} \citet{readerScaling2026} đề xuất đo mức bảo toàn
nâng cấp và cho thấy bộ nén cố định làm sụp độ co giãn theo reader. Chúng tôi dùng
độ đo này làm biến phụ thuộc của RQ3.

\begin{table*}[t]
\centering
\footnotesize
\setlength{\tabcolsep}{4pt}
\begin{tabular}{lllll}
\toprule
\textbf{Phương pháp} & \textbf{Tín hiệu giám sát} & \textbf{Cách tạo nhãn} & \textbf{Chi phí suy luận} & \textbf{Tách reader} \\
\midrule
ContextCite & log-prob đáp án & mặt nạ ngẫu nhiên + Lasso & $K$ lần gọi reader & không \\
AT2 & log-prob đáp án & mặt nạ ngẫu nhiên $\to$ đầu chú ý & 1 lượt reader & không \\
LooComp & nhãn nhị phân & leave-one-out tất định & 1 lượt encoder & có \\
Provence / XProvence & độ liên quan + gán chuỗi & nhãn từ LLM & 1 lượt encoder & có \\
RECOMP (trích xuất) & câu hữu ích cho LM & nhãn từ LM & 1 lượt encoder & có \\
EXIT & câu hữu ích & nhãn phân loại câu & 1 lượt LM 2B/câu & có \\
LongLLMLingua & perplexity có điều kiện & không huấn luyện & nhiều lượt LM 7B & có \\
CORE-RAG & EM & phần thưởng RL & sinh tự hồi quy & có \\
\midrule
\textbf{Của chúng tôi} & \textbf{F1 đầu ra} & \textbf{mặt nạ + Ridge, $\geq$3 reader} & \textbf{1 lượt encoder} & \textbf{có} \\
\bottomrule
\end{tabular}
\caption{Vị trí của phương pháp so với các công trình gần nhất.}
\label{tab:positioning}
\end{table*}

% ============================================================================
\section{Bài toán và giả thuyết}
\label{sec:hypotheses}

\subsection{Ký hiệu}

Một ví dụ gồm câu hỏi $q$, tập đáp án tham chiếu $\mathcal{A}$ và tài liệu $d$
được chia thành $C$ đoạn $(c_1,\dots,c_C)$. Với mặt nạ $m\in\{0,1\}^C$, ký hiệu
$d_m$ là ngữ cảnh chỉ giữ các đoạn có $m_i{=}1$, theo thứ tự gốc. Reader $R$ sinh câu trả
lời $R(q,d_m)$. \emph{Mức đóng góp} (utility) của mặt nạ đối với $R$ là
\begin{equation}
  u_R(m) = \mathrm{F1}\big(R(q,d_m),\mathcal{A}\big),
  \label{eq:utility}
\end{equation}
trong đó F1 tính trên token sau chuẩn hóa. Với tiếng Việt, token là âm tiết. Một
bộ nén với tỉ lệ $\rho$ chọn một mặt nạ $\hat m$ thỏa
$\mathrm{tok}(d_{\hat m})\le B=\lceil \mathrm{tok}(d)/\rho\rceil$ và được đánh giá
bằng $u_R(\hat m)$.

\subsection{Giả định làm việc}
\label{sec:assumptions}

\begin{itemize}
  \item[\textbf{A1}] \textbf{Tính cộng tính xấp xỉ.} $u_R(m)\approx\beta_0+\sum_i\beta_i m_i$.
  Giả định này chắc chắn sai khi câu hỏi cần hai đoạn \emph{cùng lúc}. Vì vậy chúng
  tôi \emph{đo} mức độ đúng của nó bằng $R^2$ kiểm định chéo của surrogate thay vì
  mặc nhiên chấp nhận (\S\ref{sec:label-quality}).
  \item[\textbf{A2}] \textbf{Mặt nạ chung.} Mọi reader nhận cùng một tập mặt nạ,
  được sinh từ hạt giống là định danh tài liệu. Nhờ vậy, khác biệt giữa nhãn của
  các reader chỉ đến từ reader.
  \item[\textbf{A3}] \textbf{Bộ nén cố định khi đánh giá.} Lựa chọn được tính một
  lần và mọi reader trả lời trên cùng lựa chọn đó. Nhờ vậy, mức bảo toàn nâng
  cấp có ý nghĩa \citep{readerScaling2026}.
\end{itemize}

\subsection{Câu hỏi nghiên cứu và giả thuyết}

Bài báo trả lời ba câu hỏi nghiên cứu. Mỗi câu gắn với đúng một giả thuyết, lần lượt ứng với ba ý
(i)--(iii) của giả thuyết trung tâm (\S\ref{sec:intro}). Mỗi giả thuyết được phát biểu kèm \emph{cách
đo} và \emph{điều sẽ xảy ra nếu bị bác bỏ}, và được kiểm định bằng một hoặc hai họ kiểm định bootstrap
cụm ghép cặp một phía, hiệu chỉnh Holm trong họ. Các họ này được cố định trước khi chạy đầy đủ và liệt
kê ở Phụ lục~\ref{app:tests}. Ký hiệu $\Delta(X,Y)=\overline{u}(X)-\overline{u}(Y)$ là chênh lệch F1
trung bình ghép cặp theo tài liệu giữa hai nhánh $X$ và $Y$, trên cùng reader, nguồn và tỉ lệ.

\paragraph{RQ1: Có chưng cất được không?}
\emph{Một bộ tỉa chưng cất từ nhãn mức đóng góp theo F1 có giữ được F1 đầu ra ở ngân sách cố định,
trong khi chỉ tốn một lượt encoder thay vì $K{\approx}64$ lần gọi reader, hay không?}

\textbf{H1.} \emph{Với mọi nguồn và $\rho\in\{4,8\}$, \arm{ours\_beta} vượt mọi bộ nén không dùng
oracle, và kém \arm{oracle\_beta} không quá 0.05 F1.}
Các bộ nén không dùng oracle gồm năm baseline đơn giản (\arm{bm25}, \arm{embed}, \arm{reranker},
\arm{lead}, \arm{random}) và bảy bộ nén đã công bố (\S\ref{sec:setup}). Chúng thuộc cùng một họ kiểm
định, vì đây là cùng một khẳng định; hiệu chỉnh Holm trên cả họ chặt hơn so với tách riêng. Đối chứng
quan trọng nhất là \arm{reranker}, chính là xương sống \texttt{bge-reranker-v2-m3} \emph{trước khi}
huấn luyện bằng nhãn mức đóng góp: thắng nó là bằng chứng rằng lợi ích đến từ \emph{nhãn}, không phải
từ kiến trúc. Vế thứ hai là một kiểm định không-kém-hơn. \arm{oracle\_beta} dùng trực tiếp nhãn Giai
đoạn A trên tài liệu kiểm tra với chi phí $K{+}1$ lần gọi reader, tức ContextCite ở cùng ngân sách:
phiên bản \emph{chưa chưng cất}. Vế này phát biểu rằng việc chưng cất làm mất ít hơn 5 điểm F1.
\emph{Nếu bác bỏ:} không thắng \arm{reranker} nghĩa là nhãn mức đóng góp không thêm thông tin so với
độ liên quan có sẵn, và luận điểm về chưng cất không đứng vững.

\paragraph{RQ2: Mức đóng góp có hơn nhãn ``đoạn chứa đáp án'' không? (phép thử sống--còn)}
\emph{Nhãn mức đóng góp có vượt trội so với giám sát bằng nhãn đoạn chứa đáp án hay không?} Đây là
giả thuyết quyết định tính mới của bài. Nếu $\beta$ chỉ tái tạo nhãn ``đoạn chứa đáp án'', việc tốn
$K$ lần gọi reader cho mỗi tài liệu huấn luyện là vô ích.

\textbf{H2.} \emph{Trên dữ liệu nhiều bước (VIMQA, HotpotQA, 2Wiki), \arm{ours\_beta} vượt
\arm{span\_sup} và \arm{oracle\_span}; trên dữ liệu một bước (UIT-ViQuAD, XQuAD-vi), hai cách giám
sát tương đương, $|\Delta(\arm{ours\_beta},\arm{span\_sup})|<0.02$.}
Cơ chế giả định cho vế nhiều bước là các đoạn cầu nối không chứa đáp án nên bị nhãn đáp án bỏ qua,
nhưng vẫn có $\beta_i>0$ vì thiếu chúng thì reader trả lời sai. Vế một bước chỉ là \emph{dự đoán},
kiểm định bằng TOST \citep{schuirmann1987tost}: khi chỉ có một đoạn kim (needle), $\beta$ được kỳ vọng
sụp về đúng đoạn đó, điều chúng tôi đo trực tiếp bằng \texttt{gold\_recall@g}; thắng ở đây là điểm
cộng, không phải yêu cầu. \emph{Quy tắc quyết định (đặt trước):} nếu vế nhiều bước thất bại, trọng tâm
bài báo chuyển sang RQ3, và đóng góp chính là phân tích tính độc lập với reader.

\paragraph{RQ3: Độc lập với reader.}
\emph{Nhãn sinh từ một reader có chuyển được sang reader khác không, và gộp nhãn có giúp không?} Với
cặp reader yếu $W$ và mạnh $S$, mức bảo toàn nâng cấp của nhánh $X$ là
\begin{equation}
  \mathrm{UR}_X(W{\to}S)=\frac{\overline{u}_S(X)-\overline{u}_W(X)}
                           {\overline{u}_S(\arm{full})-\overline{u}_W(\arm{full})}.
  \label{eq:ur}
\end{equation}
$\mathrm{UR}{=}1$ nghĩa là bộ nén giữ nguyên khoảng cách giữa hai reader, $\mathrm{UR}{=}0$ nghĩa là
xóa khoảng cách đó, và $\mathrm{UR}{<}0$ nghĩa là đảo ngược nó (Hình~\ref{fig:ur}).

\begin{figure}[t]
\centering
% Figure 3: upgrade retention (eq. ur) as a dumbbell per compressor; illustrative values.
% One hue, two shades (weak = light, strong = dark), validated as an ordinal pair on white.
\begin{tikzpicture}[font=\scriptsize, x=9.3cm, y=1cm]
  \def\xmin{0.30}
  \def\xmax{0.75}
  \def\lab{-0.215}   % left edge of the row labels, in x units relative to \xmin
  % recessive vertical grid + tick labels
  \foreach \t in {0.3,0.4,0.5,0.6,0.7} {
    \draw[gridc, line width=0.4pt] (\t-\xmin, -2.45) -- (\t-\xmin, 0.35);
    \node[text=inkM, anchor=north] at (\t-\xmin, -2.45) {\t};
  }
  \node[text=inkS, anchor=north] at (0.5*\xmax-0.5*\xmin, -2.75) {F1};
  % rows: label / weak / strong / retention
  \foreach \y/\name/\w/\s/\ur in {0/Toàn ngữ cảnh/0.40/0.70/1, -1/Bộ nén A/0.37/0.64/0.9, -2/Bộ nén B/0.35/0.41/0.2} {
    \node[anchor=east, text=inkP] at (-0.012, \y) {\name};
    \draw[axisc, line width=1.4pt, line cap=round] (\w-\xmin, \y) -- (\s-\xmin, \y);
    \fill[vzblueL, draw=white, line width=0.8pt] (\w-\xmin, \y) circle (1.35mm);
    \fill[vzblueD, draw=white, line width=0.8pt] (\s-\xmin, \y) circle (1.35mm);
    \node[anchor=west, text=inkS] at (\xmax-\xmin+0.012, \y) {$\mathrm{UR}=\ur$};
  }
  % legend
  \fill[vzblueL] (0.02, 0.6) circle (1.1mm);
  \node[anchor=west, text=inkS] at (0.032, 0.6) {reader yếu $W$};
  \fill[vzblueD] (0.27, 0.6) circle (1.1mm);
  \node[anchor=west, text=inkS] at (0.282, 0.6) {reader mạnh $S$};
\end{tikzpicture}

\caption{Minh họa mức bảo toàn nâng cấp (công thức~\ref{eq:ur}; giá trị minh họa). Bộ nén A giữ gần
trọn khoảng cách F1 giữa reader yếu và mạnh; bộ nén B xóa gần hết khoảng cách đó, nên nâng cấp reader
gần như không còn tác dụng.}
\label{fig:ur}
\end{figure}

\textbf{H3.} \emph{Nhãn gộp giữ nâng cấp reader tốt hơn nhãn một reader:
$\mathrm{UR}_{\arm{ours\_ens}}-\mathrm{UR}_{\arm{ours\_beta}}>0$ trên các cặp reader có khoảng cách F1
toàn ngữ cảnh ít nhất 0.05, kể cả các cặp có reader giữ lại.}
Các cặp có reader \emph{giữ lại} (Qwen3-32B, không tham gia sinh nhãn) tạo thành một họ kiểm định
riêng. Đây là phép thử khó nhất: nhãn gộp có tổng quát hóa sang một reader chưa từng thấy không. Mức
đồng thuận của $\beta$ giữa các reader (tương quan Spearman) được báo cáo như một số liệu mô tả.

\paragraph{Tiếng Việt là một chiều đánh giá, không phải câu hỏi nghiên cứu riêng.} Ba trong năm nguồn
là tiếng Việt, và mọi kết quả được báo cáo theo từng nguồn, nên cả ba giả thuyết đều được kiểm định
riêng trên tiếng Việt. Phép so sánh với XProvence (zero-shot, cùng họ khởi tạo
\texttt{bge-reranker-v2-m3}, nên cô lập được hiệu ứng của \emph{giám sát}) và LLMLingua-2 trên các
nguồn này nằm trong H1. Nếu XProvence đã rất mạnh trên tiếng Việt, đóng góp chuyển sang bộ đánh giá
tiếng Việt và phân tích chéo reader thay vì hiệu năng tốt nhất.

\paragraph{Những gì giả thuyết \emph{không} khẳng định.} (i) Chúng tôi không khẳng định nén tốt hơn toàn
ngữ cảnh; \arm{full} chỉ là mốc tham chiếu. (ii) Chúng tôi không khẳng định $\beta$ vượt nhãn đáp án
trên câu hỏi một bước (H2 chỉ dự đoán tương đương). (iii) Chúng tôi không khẳng định surrogate cộng tính
là đúng. A1 là xấp xỉ, và $R^2$ thấp trên dữ liệu nhiều bước được báo cáo như một phát hiện, không bị
che giấu.

% ============================================================================
\section{Phương pháp}
\label{sec:method}

Phương pháp gồm hai giai đoạn (Hình~\ref{fig:overview}). Giai đoạn A tốn kém,
chạy ngoại tuyến và chỉ trên dữ liệu huấn luyện. Giai đoạn B tạo ra bộ nén dùng
lúc suy luận.

\begin{figure*}[t]
\centering
% Figure 1: method overview (Stage A -> Stage B -> inference). \input from main.tex inside figure*.
\begin{tikzpicture}[
    font=\footnotesize,
    node distance=3.2mm and 7mm,
    box/.style={draw=#1!55, fill=white, rounded corners=2pt, align=center, text width=4.3cm,
                inner sep=3pt, line width=0.4pt, text=inkP},
    arr/.style={-{Stealth[length=3.5pt]}, line width=0.6pt, draw=inkS},
    panel/.style={draw=#1!45, fill=#1!6, rounded corners=4pt, line width=0.5pt, inner sep=2.5mm},
    ptitle/.style={font=\footnotesize\bfseries, text=inkP, align=center},
    cost/.style={font=\scriptsize\itshape, text=inkS, align=center, text width=4.3cm},
  ]
  % ---------------- A: utility attribution (offline, training data only)
  \node[box=vzblue] (a1) {Tài liệu: câu hỏi $q$, các đoạn $c_1,\dots,c_C$};
  \node[box=vzblue, below=of a1, minimum height=2.05cm] (a2) {};
  \node[anchor=north, text width=4.3cm, align=center, inner sep=0pt, text=inkP] at ($(a2.north)+(0,-1mm)$)
    {$K\!\approx\!64$ mặt nạ ngẫu nhiên $m_k$\\[-0.5mm] {\scriptsize\color{inkS}(hàng: mặt nạ; ô màu: giữ)}};
  \node[box=vzblue, below=of a2] (a3) {3 reader trả lời trên cùng các ngữ cảnh $d_{m_k}$:\\
    $u_r(m_k)=\mathrm{F1}\big(R_r(q,d_{m_k}),\mathcal{A}\big)$};
  \node[box=vzblue, below=of a3] (a4) {Ridge: $u_r(m)\approx\beta_0+\beta^{(r)\top}m$\\
    z-score trong tài liệu $\to$ nhãn $y^{(r)}$, gộp $\bar y$};
  % mask grid inside a2: 4 masks x 8 chunks, filled = chunk kept
  \foreach \row/\bits in {0/{1,0,1,1,0,1,0,0}, 1/{0,1,1,0,1,0,1,1}, 2/{1,1,0,1,0,0,1,0}, 3/{0,0,1,1,1,1,0,1}} {
    \foreach \bit [count=\col from 0] in \bits {
      \ifnum\bit=1 \def\cellstyle{vzblue}\else \def\cellstyle{gridc}\fi
      \fill[\cellstyle] ($(a2.south)+(-1.25cm+\col*3.2mm, 1.6mm+\row*2.3mm)$) rectangle ++(2.6mm,1.7mm);
    }
  }
  \node[cost, below=2mm of a4] (acost) {$K{+}1$ lần gọi reader cho mỗi tài liệu,\\ chỉ trên dữ liệu huấn luyện};
  \draw[arr] (a1) -- (a2);
  \draw[arr] (a2) -- (a3);
  \draw[arr] (a3) -- (a4);

  % ---------------- B: distillation
  \node[box=vzorange, right=of a1] (b1) {Một lượt đọc toàn bộ ngữ cảnh:\\
    \texttt{[CLS]}\,$q$\,\texttt{[SEP]}\,$c_1\dots c_C$\,\texttt{[SEP]}};
  \node[box=vzorange, below=of b1, minimum height=2.05cm] (b2) {Cross-encoder\\ {\scriptsize\texttt{bge-reranker-v2-m3}}\\
    gộp trung bình theo đoạn\\ $\to$ điểm $s_1,\dots,s_C$};
  \node[box=vzorange, below=of b2] (b3) {Mất mát so với nhãn $y$ (hoặc $\bar y$):\\
    ListNet$(s,y)+0.5\cdot$MSE$(s,y)$};
  \node[box=vzorange, below=of b3] (b4) {Chọn epoch theo NDCG@3 trên nhãn dev\\ (không gọi reader)};
  \draw[arr] (b1) -- (b2);
  \draw[arr] (b2) -- (b3);
  \draw[arr] (b3) -- (b4);
  \draw[arr] (a4.east) -- ++(3.5mm,0) |- (b3.west);   % labels feed the loss

  % ---------------- C: inference
  \node[box=vzaqua, right=of b1] (c1) {Tài liệu mới $(q, c_1,\dots,c_C)$};
  \node[box=vzaqua, below=of c1, minimum height=2.05cm] (c2) {Bộ tỉa đã huấn luyện:\\ một lượt encoder $\to$ $s_1,\dots,s_C$};
  \node[box=vzaqua, below=of c2, minimum height=2.05cm] (c3) {};
  \node[anchor=north, text width=4.3cm, align=center, inner sep=0pt, text=inkP] at ($(c3.north)+(0,-1mm)$)
    {Giữ các đoạn điểm cao nhất vừa\\ ngân sách $B$, theo thứ tự gốc\\[-0.5mm]
    {\scriptsize\color{inkS}(ô màu: đoạn được giữ)}};
  \node[box=vzaqua, below=of c3] (c4) {Reader bất kỳ, kể cả reader\\ không tham gia sinh nhãn};
  % kept (colored) / dropped (gray) chunks inside c3
  \foreach \kept [count=\col from 0] in {1,0,0,1,0,1,0,0} {
    \ifnum\kept=1 \def\cellstyle{vzaqua}\else \def\cellstyle{gridc}\fi
    \fill[\cellstyle] ($(c3.south)+(-1.25cm+\col*3.2mm, 2.2mm)$) rectangle ++(2.6mm,2.4mm);
  }
  \node[cost, below=2mm of c4] (ccost) {một lượt encoder cho mỗi tài liệu};
  \draw[arr] (c1) -- (c2);
  \draw[arr] (c2) -- (c3);
  \draw[arr] (c3) -- (c4);
  \draw[arr] (b4.east) -- ++(3.5mm,0) |- (c2.west);   % the trained pruner

  % ---------------- panels (behind) with titles
  \begin{scope}[on background layer]
    \node[panel=vzblue, fit=(a1)(a4)(acost)] (pa) {};
    \node[panel=vzorange, fit=(a1.north -| b1)(b1)(b4)(acost.south -| b1)] (pb) {};
    \node[panel=vzaqua, fit=(a1.north -| c1)(c1)(c4)(ccost)(acost.south -| c1)] (pc) {};
  \end{scope}
  \node[ptitle, above=0.5mm of pa.north] {A. Đo mức đóng góp (ngoại tuyến)};
  \node[ptitle, above=0.5mm of pb.north] {B. Chưng cất vào bộ tỉa};
  \node[ptitle, above=0.5mm of pc.north] {C. Suy luận};
\end{tikzpicture}

\caption{Tổng quan phương pháp. (A) Với mỗi tài liệu huấn luyện, ba reader trả lời trên $K$ ngữ
cảnh bị che ngẫu nhiên; hồi quy Ridge của F1 theo mặt nạ cho mức đóng góp $\beta$ của từng đoạn,
chuẩn hóa thành nhãn $y$ (một reader) hoặc $\bar y$ (gộp). (B) Nhãn được chưng cất vào một
cross-encoder chấm điểm mọi đoạn trong một lượt. (C) Lúc suy luận chỉ bộ tỉa chạy; reader chỉ được
gọi ở (A), và chỉ trên dữ liệu huấn luyện.}
\label{fig:overview}
\end{figure*}

\subsection{Giai đoạn A: đo mức đóng góp của từng đoạn}
\label{sec:method-attr}

Với mỗi tài liệu và mỗi reader, chúng tôi lấy $K=\mathrm{clamp}(C{+}1,64,256)$
mặt nạ. Trong mỗi mặt nạ, đoạn được giữ độc lập với xác suất
$p\in\{0.5,0.25\}$, xen kẽ giữa các mặt nạ, để phủ cả vùng nén vừa và nén mạnh. Hạt
giống là định danh tài liệu (A2). Reader trả lời trên từng ngữ cảnh bị che và trên
ngữ cảnh đầy đủ, rồi ta tính $u_R(m_k)$ theo công thức~(\ref{eq:utility}). Hệ số mức đóng góp là
\begin{equation}
  \begin{split}
  \hat\beta = \arg\min_{\beta_0,\beta}\;&\sum_{k=1}^{K}\big(u_R(m_k)-\beta_0-\beta^\top m_k\big)^2\\
  &+ \alpha\lVert\beta\rVert_2^2,
  \end{split}
  \label{eq:ridge}
\end{equation}
với hệ số chặn $\beta_0$ không bị phạt. Một giá trị $\alpha$ chung cho toàn bộ dữ
liệu được chọn trên lưới $\{0.1,0.3,1,3,10\}$ bằng MSE kiểm định chéo 5 phần trên các tài
liệu \emph{dev}. Nhãn giả là z-score trong tài liệu,
$y_i=(\hat\beta_i-\mathrm{mean}(\hat\beta))/\mathrm{sd}(\hat\beta)$. Tài liệu có
$\hat\beta$ hằng số (reader luôn đúng hoặc luôn sai) được đánh dấu \emph{không
mang thông tin} và bị loại khỏi huấn luyện.

Chuẩn hóa trong tài liệu làm nhãn chỉ mang thông tin \emph{thứ hạng tương đối}
giữa các đoạn của cùng tài liệu. Đây đúng là điều bộ chọn cần, và cũng loại bỏ
khác biệt về thang đo giữa các reader.

\subsection{Nhãn gộp nhiều reader}
\label{sec:method-ens}

Với các reader $r=1..3$ mà tài liệu mang thông tin, nhãn gộp là
$\bar y = z\big(\tfrac{1}{|\mathcal{R}|}\sum_{r\in\mathcal{R}} y^{(r)}\big)$. Vì
mọi reader thấy cùng mặt nạ (A2), trung bình này loại bỏ phần đóng góp riêng của
từng reader và giữ lại phần chung. Theo H3, chính phần chung này làm bộ nén ít phụ
thuộc reader hơn.

\subsection{Giai đoạn B: chưng cất vào bộ tỉa}
\label{sec:method-pruner}

Bộ tỉa là một cross-encoder đọc
\texttt{[CLS]}~$q$~\texttt{[SEP]}~$c_1\dots c_C$~\texttt{[SEP]}. Mỗi đoạn được
token hóa riêng để biết chính xác khoảng của nó. Biểu diễn ngữ cảnh hóa của các token
trong đoạn $c_i$ được gộp trung bình rồi đưa qua một lớp tuyến tính để cho điểm
$s_i$. Như vậy, điểm của mỗi đoạn được tính \emph{có điều kiện trên toàn bộ ngữ
cảnh}, không độc lập như reranker thông thường. Tài liệu dài hơn 4096 token được chia thành các cửa sổ gồm
các đoạn nguyên vẹn, mỗi cửa sổ lặp lại câu hỏi. Xương sống là
\texttt{BAAI/bge-reranker-v2-m3} \citep{chen2024bgem3}, cùng họ khởi tạo với
XProvence. Hàm mất mát là
\begin{equation}
  \begin{split}
  \mathcal{L} = &-\sum_{i}\mathrm{softmax}(y/\tau)_i\log\mathrm{softmax}(s)_i\\
  &+ \lambda\cdot\tfrac{1}{C}\sum_i (s_i-y_i)^2,
  \end{split}
  \label{eq:loss}
\end{equation}
với số hạng đầu là ListNet \citep{cao2007listnet}, $\tau=1$ và $\lambda=0.5$. Chúng
tôi chọn epoch bằng NDCG@3 trên nhãn dev, \emph{không cần gọi reader}. Nhánh đối
chứng \arm{span\_sup} dùng cùng kiến trúc, huấn luyện bằng BCE trên nhãn đoạn vàng,
và chọn epoch theo recall đoạn vàng ở ngân sách 25\%.

\subsection{Chọn đoạn trong ngân sách}
Với ngân sách $B=\lceil\mathrm{tok}(d)/\rho\rceil$, đếm bằng một bộ token hóa tham chiếu cố định, mọi
nhánh chọn theo đoạn đều dùng cùng một quy tắc: duyệt đoạn theo điểm giảm dần, giữ
những đoạn còn vừa ngân sách, rồi xuất theo thứ tự gốc. Nhờ vậy, khác biệt giữa các
nhánh chỉ nằm ở \emph{điểm số}.

\subsection{Chi phí}
Sinh nhãn cần $\sum_d (K_d+1)$ lần gọi reader cho mỗi reader. Chi phí này chỉ trả
một lần trên tập huấn luyện. Lúc suy luận, bộ tỉa cần một lượt encoder cho mỗi tài
liệu, còn \arm{oracle\_beta} (ContextCite ở cùng ngân sách) cần $K{+}1$ lần gọi
reader. Bảng~\ref{tab:cost} báo cáo cả hai chi phí.

% ============================================================================
\section{Thiết lập thực nghiệm}
\label{sec:setup}

\paragraph{Dữ liệu.} Bảng~\ref{tab:data}. Với UIT-ViQuAD \citep{nguyen2020viquad}
và XQuAD-vi \citep{artetxe2020xquad}, chúng tôi xây haystack: đoạn kim được chèn cùng
các đoạn nhiễu đến khoảng 30k ký tự, \emph{độ sâu của kim được lấy đều} để tránh
trường hợp cắt đầu/cuối thắng chỉ nhờ cách xây dữ liệu. Với VIMQA \citep{vimqa2022},
HotpotQA \citep{yang2018hotpotqa} và 2Wiki \citep{ho2020twowiki}, mỗi đoạn là một
đoạn văn có tiêu đề trong thiết lập có nhiễu gốc. Câu hỏi có/không bị loại.
XQuAD-vi và 2Wiki \emph{không bao giờ} được dùng để huấn luyện, nên chúng đo khả năng chuyển
giao chéo tập dữ liệu. Phân chia dev/test dùng \emph{hàm băm} của khóa ví dụ
(tiêu đề, đoạn văn hoặc id), không dùng hạt giống ngẫu nhiên, để tránh rò rỉ.

\begin{table}[t]
\centering
\small
\setlength{\tabcolsep}{3pt}
\footnotesize
\begin{tabular}{llllc}
\toprule
\textbf{Nguồn} & \textbf{NN} & \textbf{Bước} & \textbf{Đoạn} & \textbf{HL} \\
\midrule
UIT-ViQuAD & vi & 1 & kim + nhiễu, 30k ký tự & \checkmark \\
XQuAD-vi & vi & 1 & như trên & -- \\
VIMQA & vi & nhiều & 10 đoạn có tiêu đề & \checkmark \\
HotpotQA & en & nhiều & 10 đoạn (distractor) & \checkmark \\
2Wiki & en & nhiều & 10 đoạn có tiêu đề & -- \\
\bottomrule
\end{tabular}
\caption{Nguồn dữ liệu. NN: ngôn ngữ; HL: dùng để huấn luyện. UIT-ViQuAD: tập test là toàn bộ tập
validation chính thức (19 bài báo), tập dev là 10\% số bài báo của tập train chính thức, chọn bằng hàm băm
theo tiêu đề. XQuAD-vi, HotpotQA, 2Wiki: dev/test chia bằng hàm băm theo đoạn văn hoặc id. VIMQA dùng tập
validation/test chính thức. \todo{số tài liệu mỗi tập}}
\label{tab:data}
\end{table}

\paragraph{Reader.} Có ba reader sinh nhãn: Qwen3-8B (reader chính), Qwen3-1.7B
\citep{qwen3} và Llama-SEA-LION-v3-8B \citep{sealion2024}, một mô hình hướng tới
Đông Nam Á. Có thêm một reader \emph{giữ lại} là Qwen3-32B, chỉ dùng để đánh giá. Mọi
reader được giải mã tham lam qua vLLM.

\paragraph{Các nhánh so sánh.} \emph{Baseline đơn giản:} \arm{full} (toàn ngữ cảnh),
\arm{lead}, \arm{random}, \arm{bm25} \citep{robertson2009bm25}, \arm{embed} (bge-m3),
\arm{reranker} (bge-reranker-v2-m3 zero-shot). \emph{Bộ nén đã công bố} (bảy phương
pháp, chạy bằng checkpoint và mã gốc của tác giả):
(i) theo đoạn: Provence và XProvence \citep{chirkova2025provence,xprovence2026}, dùng điểm
reranking của chúng;
(ii) theo câu: RECOMP trích xuất \citep{xu2024recomp}, dùng checkpoint NQ cho dữ liệu một bước
và HotpotQA cho dữ liệu nhiều bước, và EXIT \citep{hwang2025exit}, dùng Gemma-2B với LoRA
gốc, cho điểm $P(\text{Yes})$ của từng câu khi biết câu hỏi và đoạn chứa câu đó;
(iii) theo token: LLMLingua \citep{jiang2023llmlingua}, LongLLMLingua
\citep{jiang2024longllmlingua} với cấu hình khuyến nghị của tác giả, và LLMLingua-2
\citep{pan2024llmlingua2}. Hai phương pháp đầu dùng Qwen2.5-7B-Instruct \citep{qwen25}
làm mô hình nhỏ thay cho Llama-2-7B mặc định, để có một mô hình đa ngôn ngữ cho tiếng Việt.
Mọi bộ nén đều bị ràng buộc cùng một ngân sách token. Các phương pháp theo câu được xếp hạng
theo điểm của chúng rồi chọn tham lam trong ngân sách, thay cho ngưỡng gốc (top-1/2 câu với
RECOMP, $P(\text{Yes})\ge0.5$ với EXIT). Các phương pháp theo token được siết tỉ lệ nén cho
tới khi vừa ngân sách. EXIT, RECOMP và Provence là mô hình tiếng Anh nên chạy zero-shot trên
tiếng Việt. Chúng tôi không đưa vào các bộ nén sinh (RECOMP abstractive, CompAct
\citep{yoon2024compact}, CORE-RAG) vì không kiểm soát được độ dài đầu ra, và CORE-RAG chưa
công bố checkpoint. Các nhánh oracle (chỉ để tham chiếu) gồm
\arm{oracle\_span}, \arm{oracle\_support} và \arm{oracle\_beta}. Các nhánh của chúng tôi
gồm \arm{ours\_beta} (một reader), \arm{ours\_ens} (gộp), cùng các biến thể ablation
\arm{ours\_logprob} và \arm{ours\_posadj}, và đối chứng \arm{span\_sup}.

\paragraph{Độ đo.} Độ đo chính là F1 token của câu trả lời (âm tiết với tiếng
Việt). Chúng tôi cũng báo cáo EM, recall đáp án, recall đoạn vàng, tỉ lệ nén thực tế,
độ trễ chọn đoạn (ms/tài liệu) và mức bảo toàn nâng cấp theo công thức~(\ref{eq:ur}).

\paragraph{Giao thức thống kê.} Mọi kết quả được báo cáo \emph{theo từng nguồn},
không gộp giữa nguồn hay ngôn ngữ, vì một lần chạy trước đã cho thấy dấu của kết quả đảo chiều giữa các
cách gộp. Khoảng tin cậy dùng bootstrap cụm, với cụm là bài báo gốc (UIT-ViQuAD) hoặc
đoạn văn (XQuAD-vi). So sánh là \emph{ghép cặp} theo tài liệu. Các họ kiểm định xác nhận ở
Bảng~\ref{tab:hyp} (Phụ lục~\ref{app:tests}) được hiệu chỉnh Holm trong từng họ. Bảng ghép cặp khám phá báo cáo
$q$ theo BH \citep{benjamini1995fdr}. Một phép kiểm định được đánh dấu
\emph{budget${\neq}$} nếu số token thực tế của hai nhánh lệch nhau quá 10\%.

\paragraph{Cài đặt.} $\rho\in\{4,8\}$, 3 epoch, 3 hạt giống huấn luyện
(\S\ref{sec:ablation}). \arm{oracle\_beta} dùng 100 tài liệu test đầu tiên của mỗi
nguồn. Siêu tham số chi tiết ở Phụ lục~\ref{app:hparams}.

% ============================================================================
\begin{figure*}[t]
\centering
\datafigure{figures/pareto.pdf}{\textwidth}
\caption{F1 theo thời gian nén trên mỗi tài liệu (thang log), reader chính, tỉ lệ nén 4$\times$, từng
nguồn riêng. Đường ngang: F1 toàn ngữ cảnh. Vòng tròn rỗng: \arm{oracle\_beta}, tức mức đóng góp đo
trực tiếp với $K{+}1$ lần gọi reader (thời gian là chi phí sinh nhãn), tính trên 100 tài liệu test đầu
tiên của mỗi nguồn; so sánh ghép cặp với \arm{ours\_beta} trên cùng các tài liệu đó là họ H1-oracle. Số
liệu từng nhánh ở Bảng~\ref{tab:main} và Bảng~\ref{tab:cost}.}
\label{fig:pareto}
\end{figure*}

\begin{figure*}[t]
\centering
\datafigure{figures/retention.pdf}{\textwidth}
\caption{Mức bảo toàn nâng cấp (công thức~\ref{eq:ur}) của \arm{ours\_beta} và \arm{ours\_ens}, tỉ lệ nén
4$\times$, với khoảng tin cậy 95\% (bootstrap cụm), cho mọi cặp reader có khoảng cách F1 toàn ngữ cảnh
ít nhất 0.05. *: cặp có reader giữ lại (không tham gia sinh nhãn). Đường dọc: 0 (xóa hết khoảng cách)
và 1 (giữ trọn).}
\label{fig:retention}
\end{figure*}

\section{Kết quả}
\label{sec:results}

\todo{Điền sau lần chạy đầy đủ (N\_TRAIN=3000). Các bảng dưới đây là khung;
số liệu lấy từ \texttt{results/eval\_test/report.json}.}

\subsection{RQ1: Chưng cất}

\begin{table*}[t]
\centering
\small
\begin{tabular}{l cc cc cc cc cc}
\toprule
& \multicolumn{2}{c}{UIT-ViQuAD} & \multicolumn{2}{c}{XQuAD-vi} & \multicolumn{2}{c}{VIMQA} & \multicolumn{2}{c}{HotpotQA} & \multicolumn{2}{c}{2Wiki} \\
\cmidrule(lr){2-3}\cmidrule(lr){4-5}\cmidrule(lr){6-7}\cmidrule(lr){8-9}\cmidrule(lr){10-11}
\textbf{Nhánh} & 4$\times$ & 8$\times$ & 4$\times$ & 8$\times$ & 4$\times$ & 8$\times$ & 4$\times$ & 8$\times$ & 4$\times$ & 8$\times$ \\
\midrule
\arm{full} & \multicolumn{2}{c}{\tbd} & \multicolumn{2}{c}{\tbd} & \multicolumn{2}{c}{\tbd} & \multicolumn{2}{c}{\tbd} & \multicolumn{2}{c}{\tbd} \\
\midrule
\arm{lead} & & & & & & & & & & \\
\arm{random} & & & & & & & & & & \\
\arm{bm25} & & & & & & & & & & \\
\arm{embed} & & & & & & & & & & \\
\arm{reranker} & & & & & & & & & & \\
\midrule
\arm{provence}$^\ddagger$ & & & & & & & & & & \\
\arm{xprovence} & & & & & & & & & & \\
\arm{recomp}$^\ddagger$ & & & & & & & & & & \\
\arm{exit}$^\ddagger$ & & & & & & & & & & \\
\arm{llmlingua} & & & & & & & & & & \\
\arm{longllmlingua} & & & & & & & & & & \\
\arm{llmlingua2} & & & & & & & & & & \\
\midrule
\arm{span\_sup} & & & & & & & & & & \\
\midrule
\arm{ours\_beta} & & & & & & & & & & \\
\arm{ours\_ens} & & & & & & & & & & \\
\midrule
\textit{\arm{oracle\_span}} & & & & & & & & & & \\
\textit{\arm{oracle\_beta}} & & & & & & & & & & \\
\bottomrule
\end{tabular}
\caption{F1 token với reader chính (Qwen3-8B) ở tỉ lệ nén 4$\times$ và 8$\times$.
Nhóm thứ nhất: baseline đơn giản; nhóm thứ hai: bộ nén đã công bố, cùng ngân sách token
(Phụ lục~\ref{app:baselines}); nhóm thứ ba: đối chứng nhãn đáp án. $^\ddagger$: mô hình chỉ
huấn luyện cho tiếng Anh, zero-shot trên các nguồn tiếng Việt. $^\dagger$: \arm{ours\_beta}
vượt trội có ý nghĩa sau hiệu chỉnh Holm trong họ H1. \arm{oracle\_beta} chỉ tính trên 100 tài liệu
test đầu tiên của mỗi nguồn (các nhánh khác: toàn bộ tập test). Kết quả của các reader khác ở Phụ lục.}
\label{tab:main}
\end{table*}

\begin{table}[t]
\centering
\small
\footnotesize
\setlength{\tabcolsep}{2pt}
\begin{tabular}{llcc}
\toprule
\textbf{Nhánh} & \textbf{Mô hình} & \textbf{Lượt} & \textbf{Thời gian} \\
\midrule
Sinh nhãn & reader 8B & $K{+}1$ & -- s \\
\arm{oracle\_beta} & reader 8B & $K{+}1$ & -- s \\
\midrule
\arm{provence} & DeBERTa-v3 & $C$ & -- ms \\
\arm{xprovence} & bge-reranker-m3 & $C$ & -- ms \\
\arm{recomp} & BERT-base & $S{+}1$ & -- ms \\
\arm{exit} & Gemma-2B & $S$ & -- ms \\
\arm{llmlingua} & Qwen2.5-7B & nhiều & -- ms \\
\arm{longllmlingua} & Qwen2.5-7B & nhiều & -- ms \\
\arm{llmlingua2} & XLM-R large & 1 & -- ms \\
\midrule
\arm{ours\_beta} & bge-reranker-m3 & $\lceil T/4096\rceil$ & -- ms \\
\bottomrule
\end{tabular}
\caption{Chi phí lúc nén cho mỗi tài liệu. Hàng đầu: đo mức đóng góp trực tiếp (sinh nhãn,
và \arm{oracle\_beta} tức ContextCite ở cùng ngân sách). Lượt: số lần chạy mô hình ($K{+}1\approx65$ với reader); $C$: số đoạn, $S$: số
câu, $T$: số token của tài liệu; bge-reranker-m3 là \texttt{bge-reranker-v2-m3}. Thời gian đo
trên một H100, gồm cả các lần siết lại tỉ lệ nén của nhánh theo token. \tbd{} Điền thời gian.}
\label{tab:cost}
\end{table}

\tbd{} Nhận xét H1: có vượt \arm{reranker} không, so với từng bộ nén đã công bố theo nguồn và
tỉ lệ (nêu rõ các ô bị đánh dấu budget${\neq}$ của nhánh theo token), và khoảng cách tới
\arm{oracle\_beta}. So sánh độ trễ ở Bảng~\ref{tab:cost}: bộ tỉa chạy
một lượt encoder, trong khi EXIT và LongLLMLingua chạy một mô hình sinh 2--7B. Hình~\ref{fig:pareto}
đặt F1 cạnh thời gian nén của mọi bộ nén.

\subsection{RQ2: Mức đóng góp so với nhãn đáp án}
\tbd{} Trình bày H2: vế nhiều bước trên ba nguồn nhiều bước và vế một bước (TOST) trên hai nguồn,
kèm \texttt{gold\_recall@g} của $\beta$ để kiểm tra hiện tượng sụp về kim.

\subsection{RQ3: Độc lập với reader}
\tbd{} Ma trận Spearman chéo reader (mô tả); mức bảo toàn nâng cấp của \arm{ours\_beta}
và \arm{ours\_ens} theo từng cặp reader, tách riêng các cặp có Qwen3-32B (H3; Hình~\ref{fig:retention}).

\subsection{Kết quả trên tiếng Việt}
\tbd{} H1--H3 riêng trên UIT-ViQuAD, XQuAD-vi và VIMQA; nhấn mạnh so sánh với \arm{xprovence} và
\arm{llmlingua2} (một phần của H1).

% ============================================================================
\section{Phân tích và ablation}
\label{sec:ablation}

\subsection{Chất lượng nhãn}
\label{sec:label-quality}
Chúng tôi đo chất lượng nhãn thay vì mặc nhiên tin vào nó. Với mỗi (reader, nguồn),
chúng tôi báo cáo: $R^2$ kiểm định chéo của surrogate cộng tính (kiểm tra A1; kỳ vọng
thấp hơn trên dữ liệu nhiều bước do tương tác giữa các đoạn), \texttt{gold\_recall@g} và MRR
của đoạn vàng theo $\beta$, tỉ lệ tài liệu mang thông tin, và tỉ lệ phương sai
nhãn được giải thích bởi một tiên nghiệm vị trí gộp (\texttt{position\_r2}). \tbd{}

\subsection{Ablation}
\begin{itemize}
  \item \textbf{F1 so với log-prob} (\arm{ours\_logprob}): phân biệt trực tiếp với AT2
  và ContextCite, vốn dùng log-prob.
  \item \textbf{Hiệu chỉnh vị trí} (\arm{ours\_posadj}): trừ tiên nghiệm vị trí khỏi nhãn để bộ
  tỉa không học thiên lệch lost-in-the-middle của reader \citep{liu2024lost}. Chúng tôi
  cũng báo cáo F1 theo ngũ phân vị độ sâu của kim (Hình~\ref{fig:depth}, Phụ lục~\ref{app:depth}).
  \item \textbf{Một reader so với gộp}, kể cả trên reader giữ lại.
  \item \textbf{Nhiễu khó so với nhiễu dễ}: nhiễu cùng bài báo cho dữ liệu một bước; kéo dài tài
  liệu nhiều bước bằng nhiễu dễ.
  \item \textbf{Hạt giống huấn luyện}: độ lệch chuẩn giữa 3 hạt giống và số giả thuyết mà
  mọi hạt giống cùng đồng ý.
\end{itemize}
\tbd{}

% ============================================================================
\section{Kết luận}
Chúng tôi đề xuất và kiểm định giả thuyết rằng mức đóng góp của từng đoạn vào F1 đầu ra
có thể được chưng cất vào một bộ tỉa một lượt, rằng nó chứa thông tin vượt quá nhãn
đáp án khi câu hỏi cần nhiều đoạn, và rằng gộp nhãn nhiều reader giúp bộ nén ít phụ thuộc
reader hơn. \tbd{} Tóm tắt kết quả theo từng giả thuyết và giả thuyết nào bị bác bỏ.

% ============================================================================
\section*{Hạn chế (Limitations)}

\begin{itemize}
  \item \textbf{Surrogate cộng tính.} Ridge trong công thức~(\ref{eq:ridge}) bỏ qua tương tác giữa các đoạn,
  trong khi câu hỏi nhiều bước về bản chất cần tương tác. Chúng tôi đo mức sai lệch
  bằng $R^2$ nhưng chưa thử surrogate có tương tác cặp.
  \item \textbf{Chi phí sinh nhãn} tỉ lệ với số reader $\times$ $K$ $\times$ số tài
  liệu huấn luyện. Chi phí này được trả một lần nhưng không nhỏ.
  \item \textbf{Chỉ chọn theo đoạn}: bộ tỉa chỉ giữ hoặc bỏ nguyên đoạn, chưa cắt hoặc
  tóm tắt bên trong đoạn. Ngân sách đạt được nên có thể thấp hơn ngân sách danh nghĩa.
  \item \textbf{F1 token} là thước đo không hoàn hảo cho câu trả lời tự do. Với tiếng
  Việt, F1 theo âm tiết có thể đánh giá quá cao các câu trả lời trùng một phần.
  \item \textbf{Phạm vi reader}: ba reader sinh nhãn đều dưới 10B tham số và reader giữ lại
  cùng họ Qwen. Khả năng chuyển sang các họ mô hình khác, ví dụ mô hình thương mại, chưa được kiểm chứng.
  \item \textbf{Haystack tổng hợp} cho dữ liệu một bước không phản ánh đầy đủ phân phối
  truy xuất thật.
  \item \textbf{Baseline}: các bộ nén sinh (RECOMP abstractive, CompAct, CORE-RAG) không được
  so sánh vì không kiểm soát được độ dài đầu ra. EXIT, RECOMP và Provence chỉ được huấn luyện
  cho tiếng Anh, nên kết quả tiếng Việt của chúng là zero-shot. Việc ép mọi bộ nén vào cùng ngân
  sách thay cho ngưỡng gốc của chúng có thể khác cấu hình tốt nhất mà tác giả báo cáo.
\end{itemize}

\section*{Cân nhắc đạo đức}
Chúng tôi dùng các bộ dữ liệu công khai theo điều khoản nghiên cứu của chúng. UIT-ViQuAD
và VIMQA yêu cầu thỏa thuận người dùng. \todo{xác nhận giấy phép trước khi nộp}
Bộ nén có thể loại bỏ bằng chứng quan trọng và dẫn tới câu trả lời sai tự tin. Vì vậy
chúng tôi không khuyến nghị dùng nó trong các lĩnh vực có rủi ro cao mà không có kiểm
tra bổ sung.

\bibliography{references}

% ============================================================================
\appendix

\section{Họ kiểm định xác nhận}
\label{app:tests}

\begin{table*}[t]
\centering
\small
\setlength{\tabcolsep}{4pt}
\begin{tabular}{L{1.4cm}L{1.9cm}L{9.2cm}L{1.9cm}}
\toprule
\textbf{Giả thuyết} & \textbf{Họ} & \textbf{Phép kiểm định (một phía, ghép cặp)} & \textbf{Phạm vi} \\
\midrule
H1 & H1 & $\arm{ours\_beta}-Y>0$, $Y\in\{$bm25, embed, reranker, lead, random$\}$ và 7 bộ nén đã công bố & mọi nguồn \\
 & H1-oracle & $\arm{ours\_beta}-\arm{oracle\_beta}>-0.05$ (không kém hơn) & mọi nguồn \\
\midrule
H2 & H2a & $\arm{ours\_beta}-Y>0$, $Y\in\{$span\_sup, oracle\_span$\}$ & nhiều bước \\
 & H2b & $|\arm{ours\_beta}-\arm{span\_sup}|<0.02$ (TOST, dự đoán) & một bước \\
\midrule
H3 & H3 & $\mathrm{UR}(\arm{ours\_ens})-\mathrm{UR}(\arm{ours\_beta})>0$, các cặp reader sinh nhãn & mọi nguồn \\
 & H3-heldout & như trên, các cặp có reader giữ lại & mọi nguồn \\
\bottomrule
\end{tabular}
\caption{Sáu họ kiểm định xác nhận cho ba giả thuyết, cố định trước khi chạy đầy đủ. Mỗi họ là một tập
kiểm định bootstrap cụm ghép cặp một phía theo (nguồn $\times$ tỉ lệ) trên reader sinh nhãn chính,
hiệu chỉnh Holm \citep{holm1979} trong họ với $\alpha=0.05$. Một giả thuyết được coi là đứng vững khi
mọi họ của nó đạt. Bảng so sánh ghép cặp đầy đủ chỉ mang tính khám phá; tương quan Spearman giữa các
reader chỉ mang tính mô tả.}
\label{tab:hyp}
\end{table*}

\section{Siêu tham số}
\label{app:hparams}
\begin{table}[h]
\centering
\small
\begin{tabular}{L{2.9cm}L{3.9cm}}
\toprule
Số mặt nạ $K$ & $\mathrm{clamp}(C{+}1,64,256)$ \\
Xác suất giữ $p$ & $\{0.5, 0.25\}$ xen kẽ \\
Lưới $\alpha$ & $\{0.1,0.3,1,3,10\}$, CV 5 phần trên dev \\
Xương sống & \texttt{BAAI/}\allowbreak\texttt{bge-reranker-v2-m3} \\
Độ dài tối đa & 4096 token/cửa sổ \\
Mất mát & ListNet ($\tau{=}1$) + $0.5\cdot$MSE \\
Số epoch & 3, chọn theo NDCG@3 dev \\
Tỉ lệ nén & 4$\times$, 8$\times$ \\
Số lần lặp bootstrap & 5000 \\
Lề H1-oracle / H2b & $0.05$ / $\pm0.02$ \\
Khoảng cách nâng cấp tối thiểu & $0.05$ \\
\bottomrule
\end{tabular}
\caption{Siêu tham số.}
\end{table}

\section{Chi tiết cài đặt baseline}
\label{app:baselines}

Bảng~\ref{tab:baselines} liệt kê checkpoint và cách mỗi bộ nén được đưa về cùng ngân sách. Mọi
nhánh nhận cùng tài liệu, cùng ngân sách $B=\lceil\mathrm{tok}(d)/\rho\rceil$ đếm bằng tokenizer
của Qwen3-8B, và lựa chọn được tính một lần rồi dùng cho mọi reader.

\begin{table*}[t]
\centering
\small
\setlength{\tabcolsep}{4pt}
\begin{tabular}{llll}
\toprule
\textbf{Phương pháp (checkpoint)} & \textbf{Đơn vị} & \textbf{Điểm dùng để xếp hạng} & \textbf{Ngưỡng gốc $\to$ ở đây} \\
\midrule
Provence & đoạn & điểm reranking & ngưỡng xác suất token $\to$ ngân sách \\
\multicolumn{4}{l}{\quad\texttt{naver/provence-reranker-debertav3-v1}} \\
XProvence & đoạn & điểm reranking & ngưỡng xác suất token $\to$ ngân sách \\
\multicolumn{4}{l}{\quad\texttt{naver/xprovence-reranker-bgem3-v1}} \\
RECOMP (trích xuất) & câu & tích vô hướng câu hỏi--câu & top-1/2 câu $\to$ ngân sách \\
\multicolumn{4}{l}{\quad\texttt{fangyuan/nq\_extractive\_compressor} (một bước)} \\
\multicolumn{4}{l}{\quad\texttt{fangyuan/hotpotqa\_extractive\_compressor} (nhiều bước)} \\
EXIT & câu & $P(\text{Yes}\mid q,\text{đoạn},\text{câu})$ & $P\ge0.5$ $\to$ ngân sách \\
\multicolumn{4}{l}{\quad\texttt{google/gemma-2b-it} + LoRA \texttt{doubleyyh/exit-gemma-2b}} \\
LLMLingua & token & perplexity & tỉ lệ $\to$ tỉ lệ siết lại \\
LongLLMLingua & token & perplexity có điều kiện theo câu hỏi & tỉ lệ $\to$ tỉ lệ siết lại \\
\multicolumn{4}{l}{\quad gói \texttt{llmlingua} 0.2.2, mô hình nhỏ \texttt{Qwen/Qwen2.5-7B-Instruct} (cho cả hai)} \\
LLMLingua-2 & token & xác suất giữ token & tỉ lệ $\to$ tỉ lệ siết lại \\
\multicolumn{4}{l}{\quad\texttt{microsoft/llmlingua-2-xlm-roberta-large-meetingbank}} \\
\bottomrule
\end{tabular}
\caption{Cài đặt các bộ nén đã công bố.}
\label{tab:baselines}
\end{table*}

\paragraph{Nhánh theo đoạn.} Provence và XProvence được dùng như reranker: mỗi đoạn nhận điểm
reranking của mô hình, rồi áp dụng cùng quy tắc chọn tham lam như các nhánh theo đoạn khác. Phần
cắt tỉa bên trong đoạn của chúng không được dùng, để mọi nhánh theo đoạn chỉ khác nhau ở điểm số.

\paragraph{Nhánh theo câu.} Câu được tách bằng cùng một bộ tách dựa trên dấu câu cho cả tiếng Việt
và tiếng Anh. Câu được xếp theo điểm của phương pháp và giữ tham lam trong ngân sách, rồi ghép lại
theo thứ tự gốc. Tiêu đề của một đoạn (dữ liệu nhiều bước) được giữ một lần khi có ít nhất một câu
của đoạn đó được chọn, và được tính vào ngân sách. RECOMP dùng checkpoint NQ cho dữ liệu một bước
và checkpoint HotpotQA cho dữ liệu nhiều bước, gộp trung bình theo token và tích vô hướng không
chuẩn hóa như mã gốc. EXIT dùng đúng mẫu prompt của tác giả, với ngữ cảnh là đoạn chứa câu. Mô hình
chạy ở bf16 thay vì lượng tử hóa 4-bit của bài gốc, với chỉ số vị trí tường minh để việc đệm trái
khi chạy theo lô không làm đổi điểm.

\paragraph{Nhánh theo token.} LLMLingua và LongLLMLingua dùng gói \texttt{llmlingua} 0.2.2. LongLLMLingua
dùng cấu hình khuyến nghị trong README (xếp hạng tài liệu theo câu hỏi, \texttt{reorder\_context=sort},
tỉ lệ nén động 0.3, \texttt{context\_budget=+100}). Câu hỏi chỉ dùng để điều kiện hóa, không được
ghép vào ngữ cảnh đầu ra. Gói này không tương thích với \texttt{transformers}~$\geq$~4.48 vì cách
nó tái sử dụng KV-cache, nên hai nhánh này chạy với \texttt{transformers} 4.46.3. Chúng tôi kiểm
tra rằng tokenizer đếm ngân sách cho cùng số token ở hai phiên bản. Vì các phương pháp theo token nén theo tỉ lệ đếm
bằng tokenizer riêng, chúng tôi siết tỉ lệ theo lượng vượt thực tế (tối đa ba lần), rồi cắt phần đuôi
nếu vẫn vượt ngân sách. Các phép so sánh mà số token thực tế của hai nhánh lệch nhau quá 10\% được
đánh dấu budget${\neq}$.

\section{F1 theo vị trí đoạn kim}
\label{app:depth}

\begin{figure}[h]
\centering
\datafigure{figures/depth.pdf}{\columnwidth}
\caption{F1 theo vị trí đoạn kim (ngũ phân vị của vị trí tương đối, 1 = đầu tài liệu) trên hai nguồn
một bước, reader chính, tỉ lệ nén 4$\times$. Một bộ nén không thiên lệch vị trí cho đường gần phẳng;
\arm{lead} chỉ giữ phần đầu tài liệu nên phụ thuộc mạnh vào vị trí.}
\label{fig:depth}
\end{figure}

\section{Tài nguyên tính toán}
Toàn bộ thí nghiệm chạy trên 4$\times$H100 80GB. Sinh nhãn chiếm phần lớn chi phí:
khoảng $3\ \text{nguồn}\times3000\ \text{tài liệu}\times65\approx0.6$M lượt sinh cho mỗi reader
(khoảng 1.8M cho ba reader). Huấn luyện mỗi bộ tỉa mất vài giờ trên một GPU. \tbd{} Thời gian thực tế.

\section{Chạy thử quy mô nhỏ}
\label{app:pilot}
Một lần chạy thử (60 tài liệu huấn luyện, 10 tài liệu kiểm tra mỗi nguồn) xác nhận
toàn bộ đường ống chạy từ đầu đến cuối. Ở quy mô đó, \arm{ours\_beta}/\arm{ours\_ens}
có F1 cao hơn \arm{lead} trong 72/80 ô (reader $\times$ nguồn $\times$ tỉ lệ), nhưng
gần như ngang \arm{embed} (35/80 ô, chênh trung bình $+0.010$). Với cỡ mẫu này,
các con số \emph{không} được dùng làm bằng chứng cho bất kỳ giả thuyết nào; chúng chỉ
cho thấy \arm{embed} là một đối thủ mạnh cần theo dõi.

\end{document}
